\documentclass[11pt]{article}

% --- Packages ---
\usepackage[margin=1in]{geometry}
\usepackage{times}
\usepackage{amsmath,amssymb}
\usepackage{booktabs}
\usepackage{hyperref}
\usepackage{enumitem}
\usepackage{graphicx}
\usepackage{xcolor}
\usepackage{natbib}
\usepackage{tabularx}
\usepackage{array}
\usepackage{caption}
\usepackage{listings}
\usepackage{fancyhdr}

\hypersetup{
    colorlinks=true,
    linkcolor=blue!70!black,
    citecolor=blue!70!black,
    urlcolor=blue!70!black
}

\lstset{
    basicstyle=\ttfamily\small,
    breaklines=true,
    frame=single,
    backgroundcolor=\color{gray!5},
    xleftmargin=1em,
    framexleftmargin=0.5em,
    columns=fullflexible,
}

\pagestyle{fancy}
\fancyhf{}
\rhead{COMPSCI 690L --- Project Plan}
\lhead{Li \& Shukla}
\cfoot{\thepage}

\title{\textbf{To Think or Not to Think: Understanding Task-Specific Reasoning in LLMs} \\ \vspace{0.3em} \large Literature Review and Project Design}
\author{
    Junyan Li \\
    UMass Amherst \\
    \texttt{junyanli@umass.edu}
    \and
    Utsav Shukla \\
    UMass Amherst \\
    \texttt{utsavshukla@umass.edu}
}
\date{April 2026}

\begin{document}
\maketitle

\begin{abstract}
This document presents the literature review and detailed project design for our study on task-specific reasoning behavior in large language models (LLMs). We survey the existing literature on chain-of-thought reasoning, overthinking phenomena, budget forcing techniques, test-time compute scaling laws, and task-specific reasoning differences. Based on this review, we present a concrete experimental design: using four model sizes from the Qwen3.5 family (0.8B, 4B, 9B, 27B), we will systematically evaluate seven diverse benchmarks under seven reasoning token budget levels (0 to 8192 tokens) using a budget forcing strategy. The goal is to characterize the accuracy--reasoning token relationship across task types and model sizes, providing practical guidance for reasoning budget allocation in LLM-based workflows.
\end{abstract}

% ============================================================
\section{Literature Review}
% ============================================================

We organize the related literature into five thematic areas that inform our study.

% ------------------------------------------------------------
\subsection{Chain-of-Thought and Thinking Modes}
% ------------------------------------------------------------

Chain-of-thought (CoT) prompting, introduced by \citet{wei2022chain}, demonstrated that providing intermediate reasoning steps as exemplars dramatically improves LLM performance on arithmetic, commonsense, and symbolic reasoning tasks. Crucially, these benefits emerge only at sufficient model scale. \citet{kojima2022large} further showed that zero-shot CoT---simply appending ``Let's think step by step''---yields similar improvements without exemplars, suggesting that LLMs have latent reasoning capabilities that can be triggered with minimal prompting. \citet{wang2023self} proposed self-consistency, sampling diverse reasoning paths and selecting the most consistent answer via majority voting, achieving +17.9\% on GSM8K over greedy CoT.

More recently, a new paradigm has emerged: \emph{native thinking modes}, where models are trained via reinforcement learning to generate internal reasoning traces. OpenAI's o1 \citep{openai2024reasoning} achieves 89th percentile on Codeforces and top-500 on AIME through RL-trained chain-of-thought. DeepSeek-R1 \citep{deepseek2025r1} demonstrates that pure RL without supervised fine-tuning on reasoning traces can incentivize reasoning capabilities, with self-reflection and verification emerging naturally. The Qwen3 family \citep{yang2025qwen3} introduces a hybrid thinking/non-thinking architecture within a single model using \texttt{<think>} tags, controlled via the \texttt{enable\_thinking} parameter. This hybrid design is directly relevant to our study, as it enables fine-grained control over reasoning behavior.

\citet{meincke2025decreasing} find that explicit CoT prompting provides diminishing returns as models become more capable, since modern reasoning models already engage in step-by-step reasoning without being asked.

% ------------------------------------------------------------
\subsection{Overthinking: When Reasoning Hurts}
% ------------------------------------------------------------

While reasoning is generally assumed to improve performance, a growing body of work demonstrates that excessive reasoning can be counterproductive. \citet{aggarwal2025optimal} introduce OptimalThinkingBench, comprising OverthinkingBench (simple queries) and UnderthinkingBench (challenging reasoning tasks), finding that independently optimizing for one metric degrades the other. \citet{su2025between} provide a systematic empirical study revealing a non-monotonic \emph{inverted-U} relationship between reasoning length and accuracy: performance improves up to a point, then declines with further reasoning. Critically, they find that LLMs tend to overthink simple problems and underthink harder ones.

\citet{hassid2025dont} demonstrate that shorter reasoning chains are up to 34.5\% more likely to yield correct answers than the longest chain, proposing \emph{short-m@k}: generating $k$ reasoning paths in parallel, stopping when the first $m$ complete, and majority voting among them. This achieves similar or superior performance using up to 40\% fewer thinking tokens.

\citet{yang2025mirage} present an empirical study showing initial performance improvements from extended thinking followed by decline, proposing a probabilistic model that reveals additional thinking increases output variance, creating an ``illusion of improved reasoning.'' They find that parallel thinking (multiple independent paths with majority vote) achieves up to 20\% higher accuracy than extended serial thinking.

\citet{zhang2025structural} identify two structural overthinking patterns: the \emph{Explorer} (excessive exploration of solution paths) and \emph{Late Landing} (excessive verification of already-correct solutions), finding that long-thinking models are 5--20$\times$ slower on simple tasks with no accuracy gains. \citet{sui2025stop} provide the first comprehensive survey on efficient reasoning, documenting the overthinking phenomenon systematically and noting that it is economically significant (OpenAI o1 costs \$60/1M generated tokens).

% ------------------------------------------------------------
\subsection{Budget Forcing and Reasoning Control}
% ------------------------------------------------------------

The question of how to \emph{control} reasoning token budgets has attracted significant attention. \citet{muennighoff2025s1} introduce the seminal \emph{budget forcing} strategy: forcefully terminating the thinking process at a predefined token threshold, or extending it by appending ``Wait'' tokens. After supervised fine-tuning on a curated 1,000-question dataset, their s1-32B model exceeds o1-preview on competition mathematics by up to 27\%. Budget forcing enables extrapolating beyond uncontrolled performance, making it the foundation for our experimental methodology.

\citet{aggarwal2025l1} propose Length Controlled Policy Optimization (LCPO), an RL method that optimizes for both accuracy and adherence to user-specified length constraints, outperforming s1's hard truncation with smoother control. \citet{li2025steering} present a fine-tuning-free approach using a lightweight predictor that models a Gamma distribution over remaining thinking length, steering generation at the token level. Under tight budgets, this significantly outperforms budget forcing. \citet{han2025token} propose the TALE framework, achieving 67\% output token cost reduction with competitive accuracy by dynamically adjusting reasoning tokens based on problem complexity.

On the commercial side, both OpenAI (via the \texttt{reasoning\_effort} parameter) and Anthropic (via the \texttt{budget\_tokens} parameter for Claude's extended thinking) have implemented API-level reasoning budget control, demonstrating industry adoption of this concept.

% ------------------------------------------------------------
\subsection{Test-Time Compute Scaling Laws}
% ------------------------------------------------------------

\citet{snell2024scaling} establish the foundational result that allocating more computation at inference time can be more efficient than scaling model parameters, with important implications for the pretraining-vs-inference compute tradeoff. \citet{yang2025tops} propose Thinking-Optimal Scaling (TOPS), revealing that scaling with longer CoTs can \emph{impair} reasoning in certain domains, and that there exists an optimal length distribution that differs across domains. This directly supports our thesis that optimal reasoning compute is task-specific.

\citet{yeo2025demystifying} study how inference compute scaling enhances reasoning, showing that SFT simplifies training but is not strictly necessary, and that error correction is inherent in base models but difficult to incentivize via RL for complex tasks. A large-scale study spanning 30B+ generated tokens across 8 open-source models \citep{artofscaling2025} finds that no single test-time scaling strategy universally dominates, and reasoning models exhibit distinct ``short-horizon'' vs.\ ``long-horizon'' trace-quality patterns.

% ------------------------------------------------------------
\subsection{Task-Specific Reasoning Differences}
% ------------------------------------------------------------

Several works have begun to investigate task-specific reasoning behavior. \citet{marjanovic2025thoughtology} systematically analyze DeepSeek-R1's reasoning patterns, discovering a ``sweet spot'' of reasoning length for each problem type---an optimal range where performance is highest, with longer chains showing substantially lower accuracy. This is the most directly related work to our study, though it examines only a single model size (DeepSeek-R1) and does not use explicit budget forcing.

\citet{zhao2025knowledge} evaluate 12 reasoning models on knowledge-intensive benchmarks, finding that increasing test-time computation does \emph{not} consistently improve accuracy for knowledge tasks, and in many cases leads to more hallucinations. This demonstrates that reasoning scaling benefits are not universal across task types. \citet{liu2025diffadapt} observe a U-shaped entropy pattern across problem difficulties, finding that 82.3\% of problems benefit from an ``Easy'' strategy, confirming that task difficulty demands fundamentally different reasoning approaches.

\citet{alomrani2025reasoning} provide a comprehensive survey introducing a two-tiered taxonomy: \emph{L1-controllability} (fixed compute budgets) and \emph{L2-adaptiveness} (dynamic scaling based on input difficulty), documenting systemic inefficiencies including underthinking on hard problems, overthinking on simple ones, and limited sensitivity to task complexity.

% ------------------------------------------------------------
\subsection{Gap and Contribution}
% ------------------------------------------------------------

While the literature establishes that (a) reasoning is not universally beneficial, (b) budget forcing can control reasoning compute, and (c) different tasks may have different optimal reasoning strategies, \textbf{no existing work has performed a systematic cross-task analysis at varying reasoning budgets using models of multiple sizes from the same family}. The Thoughtology paper \citep{marjanovic2025thoughtology} comes closest but studies only DeepSeek-R1 (a single model size) without explicit budget forcing. Our work fills this gap by using the Qwen3.5 family (0.8B to 27B) with budget forcing across seven diverse benchmarks, enabling a controlled study of how reasoning token budgets interact with task type and model scale.

% ============================================================
\section{Experimental Design}
% ============================================================

% ------------------------------------------------------------
\subsection{Models}
% ------------------------------------------------------------

We use four models from the Qwen3.5 family \citep{yang2025qwen3}, selected to span a wide range of model capacities while remaining feasible on a single NVIDIA H100 80GB GPU:

\begin{table}[h]
\centering
\caption{Model configurations. All models support multimodal inputs and native thinking modes.}
\label{tab:models}
\begin{tabular}{@{}llll@{}}
\toprule
\textbf{Model} & \textbf{Precision} & \textbf{VRAM (est.)} & \textbf{Rationale} \\
\midrule
Qwen3.5-0.8B & FP16 & $\sim$2 GB & Smallest; test if reasoning helps tiny models \\
Qwen3.5-4B & FP16 & $\sim$8 GB & Small-medium capacity \\
Qwen3.5-9B & FP16 & $\sim$18 GB & Medium capacity \\
Qwen3.5-27B & FP8 & $\sim$32 GB & Largest feasible on single H100 \\
\bottomrule
\end{tabular}
\end{table}

All Qwen3.5 models implement thinking via \texttt{<think>...</think>} tags and support vision inputs natively (unified multimodal architecture). The FP8 quantization for the 27B model leverages the H100's native FP8 Tensor Core support with negligible quality loss.

% ------------------------------------------------------------
\subsection{Benchmarks}
% ------------------------------------------------------------

We select seven benchmarks covering diverse cognitive requirements, enabling a comprehensive analysis of task-specific reasoning behavior:

\begin{table}[h]
\centering
\caption{Benchmark selection. Subset sizes are chosen to balance statistical reliability with computational feasibility.}
\label{tab:benchmarks}
\small
\begin{tabular}{@{}p{2.2cm}p{2cm}p{3.2cm}cp{2.5cm}@{}}
\toprule
\textbf{Category} & \textbf{Benchmark} & \textbf{Source} & \textbf{Size} & \textbf{Metric} \\
\midrule
Math (easy) & GSM8K & \texttt{openai/gsm8k} & 500 & Exact match \\
Math (hard) & MATH-500 & \texttt{HuggingFaceH4/MATH-500} & 500 & Exact match \\
Knowledge QA & MMLU-Pro & \texttt{TIGER-Lab/MMLU-Pro} & 500 & MC accuracy \\
Code & HumanEval+ & \texttt{evalplus} package & 164 & pass@1 \\
Summarization & XSum & \texttt{EdinburghNLP/xsum} & 500 & ROUGE + BERTScore \\
Instruction & IFEval & \texttt{google/IFEval} & $\sim$500 & Strict/loose accuracy \\
Vision + Math & MathVista & \texttt{AI4Math/MathVista} & 500 & Accuracy \\
\bottomrule
\end{tabular}
\end{table}

\textbf{Rationale for selection.} This benchmark set is designed to produce a clean narrative arc across task types:
\begin{itemize}[nosep]
    \item \textbf{Tasks expected to strongly benefit from reasoning}: MATH-500 (competition-level math), HumanEval+ (code generation requiring planning), MathVista (vision + mathematical reasoning).
    \item \textbf{Tasks with expected diminishing returns}: GSM8K (grade-school math, may saturate quickly), MMLU-Pro (knowledge QA where reasoning has moderate benefit).
    \item \textbf{Tasks where reasoning may be unnecessary or harmful}: XSum (summarization, primarily extractive), IFEval (instruction following, simple constraint checking).
\end{itemize}

For MMLU-Pro, we use stratified sampling across all 14 domains to ensure representativeness. For MathVista, we use the \texttt{testmini} split (1,000 examples with public labels) and sample 500.

% ------------------------------------------------------------
\subsection{Reasoning Budget Control}
% ------------------------------------------------------------

For each task and model, we evaluate seven reasoning token budget levels:

\begin{equation}
    t \in \{0, 64, 256, 1024, 2048, 4096, 8192\}
\end{equation}

where $t = 0$ corresponds to thinking disabled entirely (\texttt{enable\_thinking=False}) and $t > 0$ corresponds to thinking enabled with budget forcing at $t$ tokens.

\textbf{Budget forcing implementation.} We adopt a LogitsProcessor-based approach following \citet{muennighoff2025s1}:

\begin{enumerate}[nosep]
    \item Enable thinking mode (\texttt{enable\_thinking=True}).
    \item Track the token count inside \texttt{<think>...</think>} during generation.
    \item At $t - 1$ tokens: suppress all tokens except the newline character.
    \item At $t$ tokens: force the \texttt{</think>} token (token ID 248069 for Qwen3.5).
    \item Continue generation normally for the final answer.
\end{enumerate}

This ensures a graceful transition from thinking to answering, rather than abrupt truncation mid-sentence.

% ------------------------------------------------------------
\subsection{Inference Configuration}
% ------------------------------------------------------------

We serve models using vLLM with the following configuration:

\begin{lstlisting}
vllm serve Qwen/Qwen3.5-{size} --port 8000 \
  --quantization fp8 \            # for 27B only
  --reasoning-parser qwen3 \
  --max-model-len 32768
\end{lstlisting}

\textbf{Sampling parameters.} For thinking mode: temperature $= 0.6$, top-$p = 0.95$, top-$k = 20$. For non-thinking mode ($t = 0$): temperature $= 0.7$, top-$p = 0.8$, top-$k = 20$. Greedy decoding is avoided with thinking mode enabled, as it causes infinite repetition loops in Qwen3.5 models.

% ------------------------------------------------------------
\subsection{Evaluation Methodology}
% ------------------------------------------------------------

\textbf{Math tasks (GSM8K, MATH-500).} We extract the final numerical answer from the model output and compare via exact match. For MATH-500, answers are expected in \LaTeX-boxed format.

\textbf{Knowledge QA (MMLU-Pro).} Multiple-choice accuracy over 10 answer choices (10\% random baseline).

\textbf{Code (HumanEval+).} We use the \texttt{evalplus} framework for sandboxed code execution with 80$\times$ augmented test cases. We report pass@1 with a single generation per problem.

\textbf{Summarization (XSum).} We report ROUGE-1, ROUGE-2, and ROUGE-L, supplemented by BERTScore to mitigate ROUGE's limitations as a purely n-gram-based metric.

\textbf{Instruction following (IFEval).} We use the programmatic constraint checker provided with IFEval, reporting both strict accuracy (all constraints satisfied) and loose accuracy (partial credit).

\textbf{Vision + Math (MathVista).} Accuracy via exact match for numerical answers and multiple-choice questions, using the official evaluation script.

% ------------------------------------------------------------
\subsection{Analysis Plan}
% ------------------------------------------------------------

For each (model, benchmark) pair, we plot the \emph{accuracy--reasoning curve} $\text{Acc}(t)$, where $t$ is the reasoning token budget. We characterize each curve as one of:

\begin{itemize}[nosep]
    \item \textbf{Monotonically increasing}: Task benefits from more reasoning at all budget levels.
    \item \textbf{Inverted-U / saturating}: Task has an optimal reasoning budget beyond which performance plateaus or declines.
    \item \textbf{Flat}: Task does not benefit from reasoning tokens.
    \item \textbf{Declining}: Additional reasoning actively hurts performance.
\end{itemize}

We perform cross-model analysis by comparing curve shapes across model sizes to address RQ3 (does model size influence reasoning utility?). We report bootstrap 95\% confidence intervals for all accuracy estimates.

\textbf{Statistical reliability.} With 500 samples and accuracy near 50\%, the 95\% confidence interval is approximately $\pm$4.4 percentage points. With 164 samples (HumanEval), it is $\pm$7.7 points.

% ============================================================
\section{Codebase Architecture}
% ============================================================

The project codebase is organized into four modules:

\begin{itemize}[nosep]
    \item \textbf{Inference} (\texttt{src/inference/}): vLLM server management, OpenAI-compatible API client, and the budget forcing LogitsProcessor.
    \item \textbf{Data} (\texttt{src/data/}): Unified dataset loader with per-benchmark prompt formatting and answer extraction modules.
    \item \textbf{Evaluation} (\texttt{src/eval/}): Metric computation (exact match, ROUGE, BERTScore, pass@1) and task-specific judges (IFEval constraint checker, code execution sandbox).
    \item \textbf{Analysis} (\texttt{src/analysis/}): Accuracy--reasoning curve plotting, model $\times$ task $\times$ budget heatmaps, and statistical analysis.
\end{itemize}

The main experiment loop iterates over models in the outer loop (to avoid reloading weights), benchmarks in the middle loop, and reasoning budgets in the inner loop. All raw outputs (thinking tokens + final answer) are saved for post-hoc qualitative analysis.

% ============================================================
\section{Timeline}
% ============================================================

\begin{table}[h]
\centering
\caption{20-day project timeline.}
\label{tab:timeline}
\begin{tabular}{@{}clp{9cm}@{}}
\toprule
\textbf{Days} & \textbf{Phase} & \textbf{Work} \\
\midrule
1--2 & Setup & Provision H100, install vLLM, download models and datasets, verify end-to-end inference \\
3--5 & Core code & Implement API client, budget forcing, all data loaders, all evaluation metrics. Smoke test each benchmark with 5 examples on 0.8B \\
6--8 & Small models & Full experiment sweep on Qwen3.5-0.8B and Qwen3.5-4B \\
9--11 & Medium model & Full sweep on Qwen3.5-9B. Begin preliminary analysis of small model results \\
12--14 & Large model & Full sweep on Qwen3.5-27B (FP8). Slowest phase ($\sim$20--30 hrs GPU time) \\
15--17 & Analysis & Generate all accuracy--reasoning curves, heatmaps, statistical tests \\
18--19 & Write-up & Literature review, methodology, results with figures \\
20 & Buffer & Re-run failed experiments, polish, final submission \\
\bottomrule
\end{tabular}
\end{table}

\textbf{Compute estimate.} Approximately 3,500 examples $\times$ 7 budgets $\times$ 4 models $= \sim$98,000 inference calls. Estimated total GPU time: 4--5 days of continuous computation on a single H100.

% ============================================================
\section{Risks and Mitigations}
% ============================================================

\begin{table}[h]
\centering
\caption{Identified risks and mitigation strategies.}
\label{tab:risks}
\begin{tabular}{@{}p{5.5cm}p{8.5cm}@{}}
\toprule
\textbf{Risk} & \textbf{Mitigation} \\
\midrule
27B model too slow or out-of-memory & Fall back to 9B as largest model, or substitute Qwen3.5-35B-A3B (MoE with only 3B active parameters) \\
Budget forcing degrades output quality & Test early (day 3). If hard truncation fails, record natural token counts instead of forcing \\
Summarization metrics are noisy & Use both ROUGE and BERTScore. Acknowledge ROUGE limitations in write-up \\
Code evaluation requires sandboxing & Use the \texttt{evalplus} package which provides built-in sandboxed execution \\
Small models score near zero on hard benchmarks & This is itself a finding: small models cannot leverage reasoning tokens effectively \\
Data contamination & Within-model comparison controls for contamination, as all budget conditions share the same training data \\
\bottomrule
\end{tabular}
\end{table}

% ============================================================
% References
% ============================================================
\bibliographystyle{plainnat}

\begin{thebibliography}{30}

\bibitem[Aggarwal et~al.(2025a)]{aggarwal2025optimal}
P.~Aggarwal, S.~Kim, J.~Lanchantin, S.~Welleck, J.~Weston, I.~Kulikov, and S.~Saha.
\newblock Optimalthinkingbench: Evaluating over and underthinking in {LLMs}.
\newblock \emph{arXiv preprint arXiv:2508.13141}, 2025.

\bibitem[Aggarwal and Welleck(2025)]{aggarwal2025l1}
P.~Aggarwal and S.~Welleck.
\newblock {L1}: Controlling how long a reasoning model thinks with reinforcement learning.
\newblock \emph{arXiv preprint arXiv:2503.04697}, 2025.

\bibitem[Alomrani et~al.(2025)]{alomrani2025reasoning}
M.~A. Alomrani et~al.
\newblock Reasoning on a budget: A survey of adaptive and controllable test-time compute in {LLMs}.
\newblock \emph{arXiv preprint arXiv:2507.02076}, 2025.

\bibitem[DeepSeek-AI(2025)]{deepseek2025r1}
DeepSeek-AI.
\newblock {DeepSeek-R1}: Incentivizing reasoning capability in {LLMs} via reinforcement learning.
\newblock \emph{arXiv preprint arXiv:2501.12948}, 2025.

\bibitem[Han et~al.(2025)]{han2025token}
T.~Han, Z.~Wang, et~al.
\newblock Token-budget-aware {LLM} reasoning.
\newblock \emph{arXiv preprint arXiv:2412.18547}, 2025.

\bibitem[Hassid et~al.(2025)]{hassid2025dont}
M.~Hassid et~al.
\newblock Don't overthink it. Preferring shorter thinking chains for improved {LLM} reasoning.
\newblock \emph{arXiv preprint arXiv:2505.17813}, 2025.

\bibitem[Kojima et~al.(2022)]{kojima2022large}
T.~Kojima, S.~S. Gu, M.~Reid, Y.~Matsuo, and Y.~Iwasawa.
\newblock Large language models are zero-shot reasoners.
\newblock In \emph{NeurIPS}, 2022.

\bibitem[Li et~al.(2025)]{li2025steering}
J.~Li, W.~Zhao, Y.~Zhang, and C.~Gan.
\newblock Steering {LLM} thinking with budget guidance.
\newblock \emph{arXiv preprint arXiv:2506.13752}, 2025.

\bibitem[Liu et~al.(2025)]{liu2025diffadapt}
X.~Liu et~al.
\newblock {DiffAdapt}: Difficulty-adaptive reasoning for token-efficient {LLM} inference.
\newblock \emph{arXiv preprint arXiv:2510.19669}, 2025.

\bibitem[Marjanovic et~al.(2025)]{marjanovic2025thoughtology}
S.~V. Marjanovic et~al.
\newblock {DeepSeek-R1} thoughtology: Let's think about {LLM} reasoning.
\newblock \emph{arXiv preprint arXiv:2504.07128}, 2025.

\bibitem[Meincke et~al.(2025)]{meincke2025decreasing}
L.~Meincke, E.~R. Mollick, L.~Mollick, and D.~Shapiro.
\newblock Prompting science report 2: The decreasing value of chain of thought in prompting.
\newblock \emph{arXiv preprint arXiv:2506.07142}, 2025.

\bibitem[Muennighoff et~al.(2025)]{muennighoff2025s1}
N.~Muennighoff et~al.
\newblock s1: Simple test-time scaling.
\newblock \emph{arXiv preprint arXiv:2501.19393}, 2025.

\bibitem[OpenAI(2024)]{openai2024reasoning}
OpenAI.
\newblock Learning to reason with {LLMs}.
\newblock \url{https://openai.com/index/learning-to-reason-with-llms/}, 2024.

\bibitem[Snell et~al.(2024)]{snell2024scaling}
C.~Snell, J.~Lee, K.~Xu, and A.~Kumar.
\newblock Scaling {LLM} test-time compute optimally can be more effective than scaling model parameters.
\newblock \emph{arXiv preprint arXiv:2408.03314}, 2024.

\bibitem[Su et~al.(2025)]{su2025between}
J.~Su, J.~Healey, P.~Nakov, and C.~Cardie.
\newblock Between underthinking and overthinking: An empirical study of reasoning length and correctness in {LLMs}.
\newblock \emph{arXiv preprint arXiv:2505.00127}, 2025.

\bibitem[Sui et~al.(2025)]{sui2025stop}
Y.~Sui et~al.
\newblock Stop overthinking: A survey on efficient reasoning for large language models.
\newblock \emph{arXiv preprint arXiv:2503.16419}, 2025.

\bibitem[Wang et~al.(2023)]{wang2023self}
X.~Wang, J.~Wei, D.~Schuurmans, Q.~Le, E.~Chi, S.~Narang, A.~Chowdhery, and D.~Zhou.
\newblock Self-consistency improves chain of thought reasoning in language models.
\newblock In \emph{ICLR}, 2023.

\bibitem[Wei et~al.(2022)]{wei2022chain}
J.~Wei, X.~Wang, D.~Schuurmans, M.~Bosma, B.~Ichter, F.~Xia, E.~Chi, Q.~Le, and D.~Zhou.
\newblock Chain-of-thought prompting elicits reasoning in large language models.
\newblock In \emph{NeurIPS}, 2022.

\bibitem[Yang et~al.(2025a)]{yang2025qwen3}
A.~Yang, A.~Li, B.~Yang, B.~Zhang, et~al.
\newblock Qwen3 technical report.
\newblock \emph{arXiv preprint arXiv:2505.09388}, 2025.

\bibitem[Yang et~al.(2025b)]{yang2025tops}
W.~Yang et~al.
\newblock Towards thinking-optimal scaling of test-time compute for {LLM} reasoning.
\newblock \emph{arXiv preprint arXiv:2502.18080}, 2025.

\bibitem[Yang et~al.(2025c)]{yang2025mirage}
Y.~Yang et~al.
\newblock Does thinking more always help? Mirage of test-time scaling in reasoning models.
\newblock \emph{arXiv preprint arXiv:2506.04210}, 2025.

\bibitem[Yeo et~al.(2025)]{yeo2025demystifying}
E.~Yeo, Y.~Tong, M.~Niu, G.~Neubig, and X.~Yue.
\newblock Demystifying long chain-of-thought reasoning in {LLMs}.
\newblock \emph{arXiv preprint arXiv:2502.03373}, 2025.

\bibitem[Zhang et~al.(2025)]{zhang2025structural}
X.~F. Zhang et~al.
\newblock Do {LLMs} really need 10+ thoughts for ``find the time 1000 days later''? Towards structural understanding of {LLM} overthinking.
\newblock \emph{arXiv preprint arXiv:2510.07880}, 2025.

\bibitem[Zhao et~al.(2025)]{zhao2025knowledge}
J.~X. Zhao, B.~Hooi, and S.-K. Ng.
\newblock Test-time scaling in reasoning models is not effective for knowledge-intensive tasks yet.
\newblock \emph{arXiv preprint arXiv:2509.06861}, 2025.

\end{thebibliography}

\end{document}
